\documentclass[FGBook.tex]{subfiles}

\begin{document}

\drama{Bílá nemoc}{KAREL ČAPEK}

\noindent\begin{wrapfigure}{r}{0.35\textwidth}\footnotesize
\subwection{Ukázka z díla:}
\setlength{\parindent}{3pt}
MARŠÁL: Nemohu. Copak vám to mám vysvětlovat jako malému dítěti? Myslíte, že válka nebo mír záleží na mé vůli? Musím se řídit tím, co je v zájmu státu...\\
DR. GALÉN: Vaše Excelence, já jsem lékař. Jako lékař... musím být proti střílení a yperitu, že? A proti všemu, co zabíjí lidi. Proto léčím jen chudé. Ti aspoň nechtějí válku.
\end{wrapfigure}

\wection{LITERÁRNÍ TEORIE}

\subwection{Literární druh a žánr:}
\noindent drama, protifašistická tragédie (alegorické drama s prvky varování)

\subwection{Literární směr:}
\noindent meziválečná česká literatura -- demokratický proud, pragmatismus, expresionistické prvky, angažovaná literatura s prvky utopismu a kritiky společnosti.

\subwection{Jazyk a slovní zásoba:}
\noindent Spisovná čeština s odbornými latinskými názvy nemocí, anglickými a německými frázemi (odkaz na fašismus). Dialogy jsou přirozené, míra hovorovosti závisí na společenské vrstvě postavy. Krátké scénické poznámky, metafory a ironie.

\subwection{Figury:}
\noindent Dialogy s gradací, ironie, opakování, přerušené věty (u Galéna), frázovitost (u Maršála).

\subwection{Tropy:}
\noindent Alegorie (\texthl{bílá nemoc = fašismus/nenávist}), metafora, symbolismus (bílá skvrna jako morální úpadek), kontrast (Galén vs. Maršál).

\subwection{Stylistická charakteristika textu:}
\noindent Drama o třech aktech a 14 obrazech. Chronologický postup s paralelními dialogy nemocných a rodinnými scénami (dokumentují dopad nemoci na společnost). Dominují dialogy, které odhalují charaktery a ideologie. Civilní, nepatetický styl s tragickým vyzněním.

\subwection{Postavy:}
\noindent 
\textbfhl{DR. GALÉN (Doktor Dětina) --} mladý, laskavý, naivní, ale statečný lékař a pacifista. Objeví lék proti bílé nemoci, léčí však jen chudé a bezmocné, aby přinutil mocné k míru. Symbol humanismu a individuálního svědomí.\\
\textbfhl{MARŠÁL --} bezcitný diktátor (alegorie Hitlera/Mussoliniho), touží po moci a válce. Zbabělý, frázovitý, ochotný obětovat národ. Nakonec onemocní bílou nemocí.\\
\textbfhl{BARON KRÜG --} majitel zbrojovek, Maršálův přítel, pragmatický válečný profiteer. Onemocní a ze zoufalství se zastřelí.\\
\textbfhl{DVORNÍ RADA SIGELIUS --} ředitel nemocnice, panovačný, ambiciózní lékař toužící po slávě. Nepřiznává Galénovy zásluhy.\\
\textbfhl{SIGELIUSOVA DCERA / MLADÝ KRÜG --} představitelé mladé generace, ovlivnění propagandou, ale schopni změny pod tlakem.

\subwection{Děj:}
\noindent V blíže neurčené zemi (alegorie fašistické Evropy) propukne epidemie bílé nemoci (forma lepry – bílé skvrny, rozpad těla). Postihuje především starší lidi; mladí ji tajně vítají (uvolní místa). Dr. Galén objeví lék, ale léčí pouze chudé – chce tak donutit mocné k uzavření míru. Baron Krüg onemocní a žádá o léčbu; Galén požaduje zastavení zbrojení. Krüg se ze zoufalství zastřelí. Nakonec onemocní i Maršál. Pod tlakem dcery a mladého Krüga souhlasí s mírem a volá Galéna. Cestou do paláce však dav (ovlivněný propagandou) Galéna ušlape i s jeho lékem. Válka vypukne a nelze ji zastavit.

\subwection{Kompozice, prostor a čas:}
\noindent Tři akty, 14 obrazů + předmluva (autorovo varování). Chronologický postup s paralelními scénami (nemocní, rodiny). Prostor a čas neurčeny přesně – blíže neurčená země v předvečer války (alegorie 30. let 20. století).

\subwection{Význam sdělení:}
\noindent Silné varování před fašismem, nacismem a válečným šílenstvím. Bílá nemoc symbolizuje morální nákazu společnosti – slepé následování davu, nenávist a ztrátu humanity. Konflikt jedinec (Galén – svědomí) vs. moc/dav. Otázka: uvědomí si národ svou „nemoc“, nebo bude pokračovat ve válce až do záhuby? Čapek kritizuje i pasivitu a oportunismus.

\vspace*{2em}
\wection{LITERÁRNÍ HISTORIE}

\subwection{Politická situace:}
\noindent Premiéra 29. 1. 1937 (Stavovské divadlo, Praha). Reakce na nástup Hitlera (1933), italskou invazi do Habeše, španělskou občanskou válku a hrozbu 2. světové války. Protifašistický charakter díla vedl k útokům nacistů na Čapka (plánované zatčení gestapem).

\subwection{Kontext literárního vývoje:}
\noindent Meziválečná česká literatura – angažovaná, demokratická próza a drama. Čapek spojuje pragmatismus, žurnalistiku a filozofické prvky. Varování před zneužitím techniky a ideologií (jako v R.U.R. nebo Válce s mloky).

\subwection{Karel Čapek {\ssmall -- život autora:}}
\noindent (1890 Malé Svatoňovice – 1938 Praha). Český spisovatel, novinář, dramatik a filozof. Doktorát z UK, redaktor Lidových novin, dramaturg Vinohradského divadla. Spolu s bratrem Josefem členem Pátečníků, přítel T. G. Masaryka (Hovory s T. G. Masarykem). Měl Bechtěrevovu nemoc. Zemřel na zápal plic krátce před Mnichovem. Bratr Josef zahynul v koncentráku.

\subwection{Další autorova tvorba:}
\noindent Mezi nejznámější díla patří \textsc{R.U.R.}, \textsc{Válka s mloky}, \textsc{Věc Makropulos}, \textsc{Loupežník}, \textsc{Matka} a fejetony.

\subwection{Aktuálnost tématu a zpracování tématu:}
\noindent Varování před totalitními ideologiemi, populismem, davovou psychózou a zneužitím moci zůstává nadčasové (i v kontextu moderních pandemií a politických krizí).

\vspace*{2em}
\wection{OSOBNÍ NÁZOR}
\noindent Bílá nemoc je geniální alegorie, která i po téměř 90 letech působí naléhavě. Galén jako nepatetický hrdina všedního dne proti Maršálovi-diktátorovi ukazuje, jak snadno se společnost „nakazí“ nenávistí. Nejsilnější je tragický konec – dav ušlape jediného člověka, který mohl zastavit válku. Čapek varuje bez zbytečného patosu, skrz dialogy a symboliku. Rozhodně patří k vrcholům českého dramatu a povinné četbě pro pochopení 30. let.

\vfill

\noindent\begin{minipage}{\textwidth}
    {\textcolor{\wpagecolor}{\rule{\linewidth}{0.4pt}}
    \footnotesize
    \textbfhl{Alegorie --} zobrazení abstraktního jevu konkrétním obrazem (bílá nemoc = fašismus); 
    \textbfhl{Protifašistická literatura --} angažovaná tvorba varující před totalitními režimy; 
    \textbfhl{Tragédie --} dramatický žánr s tragickým koncem hrdiny a morálním ponaučením.
    }
\end{minipage}

\newpage

\end{document}